% !TeX root = RJwrapper.tex
\title{Simulating Time-to-Event Data Coupled to ODE Systems in R with simtte}


\author{by Carlos Serra Traynor}

\maketitle

\abstract{%
Realistic clinical trial simulation increasingly requires time to event outcomes whose hazard is linked to an underlying pharmacological mechanism rather than specified only through a closed form empirical model. Existing R packages for survival simulation support parametric hazards and flexible spline based models, but they do not natively couple the event process to the numerical solution of an ordinary differential equation (ODE) system such as a pharmacokinetic/pharmacodynamic (PK/PD) model. The simtte package addresses this gap by augmenting any mrgsolve model with a survival compartment satisfying dS/dt = -h(t)S(t), where the hazard h(t) may depend on any model state or derived quantity, including drug concentration, biomarker or mediator dynamics, tumour burden, or other mechanistically relevant variables. The survival process is integrated jointly with the pharmacological states using the same compiled solver, and event times are generated from the resulting survival trajectory by inverse transform sampling. The package provides a common interface for Weibull models, M spline baseline hazards, and fully user defined ODE coupled hazards, with a consistent output structure across these settings. Administrative censoring is handled through the simulated follow up horizon, and independent random censoring can be incorporated subsequently using standard R operations. We describe the package design, detail the joint integration and event time generation procedures, and illustrate the workflow through examples ranging from a simple Weibull model to a one compartment indirect response PK/PD model in which treatment suppresses a mediator that drives event risk.
}

\renewcommand{\thesection}{\arabic{section}}

\section{Introduction}\label{introduction}

Clinical trial simulation is a core component of contemporary drug development and quantitative decision-making. Regulatory agencies and pharmaceutical sponsors use simulation to inform sample size determination, evaluate operating characteristics, and support dose and design selection before initiation of patient enrollment \citep{ich_e9}. In many oncology and cardiovascular applications \citep{johnson2024longitudinal}, these evaluations depend critically on a time-to-event (TTE) model, which characterizes the distribution and timing of clinically relevant endpoints such as disease progression, death, or first hospitalization.

Standard practice in time-to-event simulation is to specify a parametric hazard model, most commonly the Weibull distribution, estimate or calibrate its parameters from historical data, and generate event times by inverse transform sampling \citep{bender2005}. This approach is well established and well supported in R through packages such as \CRANpkg{simsurv} \citep{simsurv}, \CRANpkg{survsim} \citep{survsim}, and \CRANpkg{simstudy} \citep{simstudy}. Its limitations become evident when the treatment mechanism is expected to influence the hazard function itself. For example, a drug that alters tumour dynamics through an indirect pharmacodynamic effect does not naturally imply a Weibull hazard; rather, the event risk arises from the interaction of exposure--response relationships, mediator turnover, disease progression dynamics, and between-subject variability in PK and PD parameters. In such settings, when the objective is to evaluate the impact of dose, regimen, or patient selection on survival outcomes, an empirically specified hazard is only a surrogate for the underlying mechanism, not a mechanistic representation of it.

Pharmacometric simulation platforms such as \CRANpkg{mrgsolve} \citep{mrgsolve} provide efficient large-scale simulation of ODE-based PK/PD models using compiled code. In practice, pharmacometricians routinely use \CRANpkg{mrgsolve} to generate concentration--time profiles and pharmacodynamic response trajectories for virtual patient populations, while accounting for population variability through random effects and other model components. What has been lacking, however, is a direct link between these continuous-time mechanistic trajectories and discrete time-to-event outcomes: when an ODE model defines the time course of a hazard function, event times still need to be generated in a way that is numerically consistent with the underlying mechanistic system.

The \CRANpkg{simtte} package closes this gap. Its central contribution is a general framework for integrating the survival function as an additional compartment within an \CRANpkg{mrgsolve} ODE model, such that the survival probability \(S(t)\) is integrated simultaneously with the pharmacological state variables using the same compiled solver. Event times are subsequently generated by applying inverse transform sampling to the simulated survival trajectory. The method is applicable to any \CRANpkg{mrgsolve} model: users specify a standard \CRANpkg{mrgsolve} model, add a survival compartment and the corresponding differential equation, and pass the resulting simulation output directly to \texttt{sim\_tte\_df()}. For users who do not require a custom PK/PD model, \CRANpkg{simtte} also provides built-in Weibull and M-spline implementations through the higher-level interface \texttt{sim\_tte()}, which handles model specification and simulation internally.

The remainder of this paper is organized as follows. Section \ref{background} reviews related R packages for time-to-event simulation and situates \CRANpkg{simtte} within the broader \CRANpkg{mrgsolve} and pharmacometric modeling ecosystem. Section \ref{package-design} describes the package architecture, the joint integration of survival and pharmacological dynamics, and the inverse transform sampling procedure used to generate event times. Section \ref{examples} presents three worked examples, progressing from a simple Weibull model to a fully mechanistic PK/PD-driven hazard specification. Section \ref{censoring} addresses administrative and stochastic censoring mechanisms. Finally, Section \ref{discussion} considers the broader implications of this approach and outlines potential directions for future development.

\section{Background}\label{background}

\subsection{Related R packages}\label{related-r-packages}

Several R packages support the simulation of time-to-event data, each with a distinct modeling scope and implementation strategy. Here, we briefly review commonly used packages and contrast their design with that of \CRANpkg{simtte}.

\CRANpkg{simsurv} \citep{simsurv} generates survival times from a user-defined log-hazard function and supports a broad range of parametric specifications, including competing risks and time-varying covariates. Its design is highly flexible when the hazard can be written directly as a closed-form function and evaluated within an R-based numerical root-finding procedure. However, it does not natively support hazards that depend on the evolving solution of an ODE system, because the survival process is not jointly integrated with a mechanistic model.

\CRANpkg{survsim} \citep{survsim} emphasizes simulation of multi-event and recurrent-event data and provides parametric models such as Weibull and log-normal formulations. It is well suited to more complex event structures, including illness--death settings, but it does not provide flexible spline-based baseline hazards or a mechanism for coupling event risk to ODE-driven PK/PD dynamics.

\CRANpkg{simstudy} \citep{simstudy} is a general-purpose clinical trial simulation framework in which survival outcomes constitute one component of a broader data-generation toolkit. Its survival functionality supports exponential and Weibull models, with treatment effects typically expressed through hazard ratios. As with \CRANpkg{simsurv} and \CRANpkg{survsim}, the hazard is specified as a closed-form function of time and covariates rather than as a quantity derived from the numerical solution of a differential equation system.

Across these packages, the hazard function is specified as an explicit R function of time and covariates. None of them support hazards defined by the numerical solution of an ODE system. In mechanistic applications, for example when event risk depends on drug exposure, biomarker dynamics, or disease progression, the ODE model must therefore be solved externally and its trajectory interpolated within the hazard function. This indirect approach separates the event process from the underlying mechanistic model and adds an additional numerical step, rather than treating survival as an integrated component of the same dynamic system.

\subsection{The mrgsolve ecosystem}\label{the-mrgsolve-ecosystem}

\CRANpkg{mrgsolve} \citep{mrgsolve} is an R package for simulation from ODE based models, developed and maintained by Metrum Research Group. It compiles model specifications written in a C++ like domain specific language into shared libraries, provides a high level R interface for simulation in large virtual populations, and integrates naturally with standard R data structures and packages such as \CRANpkg{dplyr}. The package is widely used in pharmacometrics for population PK and PK/PD modeling.

Within this ecosystem, \CRANpkg{simtte} provides a layer for mechanistic time to event simulation. Rather than introducing a separate simulation framework, it extends the existing \CRANpkg{mrgsolve} workflow by adding the components required for TTE generation, including the survival compartment, the inverse transform sampling procedure, and a unified R interface. As a result, users already familiar with \CRANpkg{mrgsolve} model code can specify bespoke PK/PD hazard models and generate event times directly with \texttt{sim\_tte\_df()}, without adopting a new modeling approach.

\section{Package design}\label{package-design}

\subsection{Survival function co-integration}\label{survival-function-co-integration}

The principal methodological contribution of \CRANpkg{simtte} is the formulation of the survival function as an additional compartment within the ODE system. For a hazard function \(h(t)\), the survival probability satisfies the forward Kolmogorov equation:

\[\frac{dS(t)}{dt} = -h(t)\, S(t), \quad S(0) = 1.\]

In \CRANpkg{mrgsolve} model code, this formulation is implemented by adding a compartment to represent the survival probability. In the package examples, this compartment is named \texttt{p11}, although the name itself is arbitrary. The compartment is initialized at 1, corresponding to \(S(0) = 1\), and its evolution is specified through the ODE:

\begin{verbatim}
$INIT p11 = 1

$ODE
dxdt_p11 = -p11 * HAZ;
\end{verbatim}

Here, \texttt{HAZ} denotes the instantaneous hazard evaluated at the current solver time. It may be defined as any expression involving other model states or derived quantities, such as drug concentration, receptor occupancy, mediator level, tumour size, or combinations of these variables. The ODE solver integrates \texttt{p11} jointly with the remaining system states in a single numerical pass, so that the survival trajectory remains consistent with the underlying pharmacological dynamics. Under this formulation, no separate numerical integration of the cumulative hazard is needed.

The built in Weibull model distributed with \CRANpkg{simtte} (\texttt{inst/models/weibull.cpp}) departs from this ODE co-integration pattern. For shape values below 1 the instantaneous hazard \(h(t) = \eta\,\gamma\,t^{\gamma-1}\), with \(\eta = \exp(\mu + \text{lp})\), diverges as \(t \to 0^+\), which is numerically unstable for an ODE solver evaluating derivatives near \(t = 0\). Because the Weibull survival function is available in closed form, \(S(t)\) is instead evaluated directly on the reported output grid:

\begin{verbatim}
$TABLE
double p11 = exp(-exp((mu + lp) + shape * log(TIME)));
\end{verbatim}

with \(t = 0\) handled as a separate case, since \(S(0) = 1\) exactly regardless of shape. This closed-form evaluation is exact to floating point precision for every shape \(> 0\) and every reported time, including near-zero and large-parameter regimes that would otherwise risk numerical overflow; agreement to floating point precision was confirmed directly against two independent references, a hand-derived closed form and R's own \texttt{stats::pweibull()}, across a broad grid of shape, intercept, and linear-predictor values (package validation scripts, \texttt{inst/validation/}). It is specific to the built in Weibull model, whose survival function has a closed form; it is not a general continuous-time event-root solver. Under this parameterization, \(\mu\) acts multiplicatively on the hazard, so it is most naturally interpreted as a rate like term on the hazard scale rather than the scale parameter used in some alternative Weibull parameterizations.

The M spline model (\texttt{inst/models/ms.cpp}), by contrast, retains the ODE co-integration pattern, defining the hazard using the baseline hazard value \texttt{basehaz}, supplied as a time varying input:

\begin{verbatim}
$ODE
dxdt_p11 = -p11 * basehaz * eta;
\end{verbatim}

\texttt{basehaz} is supplied only at the elements of the \texttt{time} grid; between two consecutive supplied times, the hazard value at the \emph{earlier} time applies (a last-observation-carried-forward convention), so the hazard used by the solver is piecewise constant on the reported grid rather than a continuous re-evaluation of the underlying M spline basis at arbitrary times. The M spline model therefore follows the same general pattern as a user defined mechanistic hazard model: the expression \texttt{dxdt\_p11\ =\ -p11\ *\ HAZ} provides the template, with the specific form of \texttt{HAZ} determined by the model. User defined mechanistic hazard models, which in general lack a closed form for \(S(t)\), should continue to use this ODE co-integration template rather than the Weibull model's closed-form evaluation.

\subsection{Inverse transform sampling}\label{inverse-transform-sampling}

Given the simulated trajectory of \(S(t)\) from the ODE solution, event times are generated by inverse transform sampling \citep{bender2005}. If
\(T\) denotes the event time with survival function \(S(t)=P(T>t)\), the standard inverse transform construction implies that

\[T = \inf\{t : S(t) \leq U\}, \quad U \sim \text{Uniform}(0,1).\]

Accordingly, for each simulated subject the algorithm draws a uniform random variable \(U\) and identifies the earliest time at which the survival probability falls below that threshold. This value is taken as the simulated event time. If \(S(t)\) does not fall below \(U\) within the modeled follow up period, the subject is administratively censored at the end of follow up.

In the current R implementation, \texttt{sim\_tte\_df()} (via \texttt{.simulate\_survival\_id()}) applies this construction to the discrete survival trajectory returned by \CRANpkg{mrgsolve} and locates the first observation for which \texttt{p11\ \textless{}=\ U}. The set of times at which the trajectory is reported, and hence the resolution at which the crossing can be identified, is controlled directly by the user supplied \texttt{time} argument, which \CRANpkg{simtte} passes to \CRANpkg{mrgsolve} through its \texttt{tgrid} mechanism rather than leaving it to \CRANpkg{mrgsolve}'s own default output schedule. Under the default \texttt{event\_time\_method\ =\ "grid"}, the estimated event time is exactly one of these reported times, so its precision is tied to the spacing of \texttt{time}, at the usual cost of finer grids requiring more computation. \CRANpkg{simtte} also provides an opt-in refinement of this resolution, described next.

\subsection{Event-time interpolation}\label{event-time-interpolation}

Setting \texttt{event\_time\_method\ =\ "log\_survival"} refines a crossing that falls strictly after the first reported observation, without changing which subjects are classified as events versus censored, and without changing the reported time for censoring or for a crossing already present at the first reported observation. Between the two reported points \((t_i, S_i)\) and \((t_{i+1}, S_{i+1})\) surrounding the crossing, the method linearly interpolates cumulative hazard \(H(t) = -\log S(t)\),

\[H(t) = H_i + \frac{H_{i+1} - H_i}{t_{i+1} - t_i}(t - t_i), \qquad H_i = -\log S(t_i),\]

and solves for the time at which \(H(t)\) equals \(-\log U\). This is equivalent to assuming the hazard is constant over \([t_i, t_{i+1})\). For the M spline model, whose hazard is genuinely piecewise constant on the reported grid by construction (Section \ref{package-design}), this recovers the event time implied by that discretized hazard exactly. For a continuously varying hazard, such as the Weibull model with \(\gamma \neq 1\) or a custom mechanistic model, it is an approximation whose accuracy improves as \texttt{time} is refined; it is exact for the Weibull model only in the special case \(\gamma = 1\), the one setting in which the Weibull hazard is itself constant. \texttt{event\_time\_method\ =\ "log\_survival"} therefore provides sub-grid event-time resolution without altering the event process itself, but it is not a general continuous-time root solver for an arbitrary reported trajectory.

This behavior has been checked empirically against an independent closed-form reference (package validation scripts, \texttt{inst/validation/}): as \texttt{time} is refined, the error of \texttt{event\_time\_method\ =\ "grid"} decreases at approximately the rate expected of a nearest-grid-point lookup, while \texttt{"log\_survival"} converges faster and matches the reference to floating point precision, at every grid spacing tested, for the exponential (\(\gamma = 1\)) case.

\subsection{Software interface}\label{software-interface}

The package exposes two primary simulation interfaces, together with one diagnostic utility and one example dataset. The main user facing functions are \texttt{sim\_tte()}, which supports the built in Weibull and M spline models, and \texttt{sim\_tte\_df()}, which operates on user generated \CRANpkg{mrgsolve} output containing a survival probability trajectory. This separation reflects the two intended use cases of the package: direct simulation from standard hazard models and simulation from custom mechanistic models in which the hazard is defined within an ODE system.

Both \texttt{sim\_tte()} and \texttt{sim\_tte\_df()} additionally accept an \texttt{event\_time\_method} argument (\texttt{"grid"}, the default, or \texttt{"log\_survival"}) controlling how the event time is resolved from the reported survival trajectory; see the Event-time interpolation section above.

A diagnostic function, \texttt{explore\_pi\_tq\_surv()}, is included to examine how the prognostic index relates to survival probability under a given hazard specification. The package also provides the example dataset \texttt{ms\_data}, which contains the components required to reproduce the M spline examples used in this article.

Across these interfaces, the output format is consistent: simulations are returned as a data frame containing the simulated event or censoring time, event indicator, subject identifier, and prognostic index. This common structure supports routine downstream analysis with standard survival analysis workflows in R.

\section{Examples}\label{examples}

\subsection{Weibull parametric model}\label{weibull-parametric-model}

A simple use case for \CRANpkg{simtte} is simulation from a Weibull proportional hazards model. The following code generates data for 50 subjects with individual log hazard ratios sampled from a standard normal distribution, a Weibull shape parameter of 1.1, and administrative censoring at time 100.

\begin{verbatim}
library(simtte)
set.seed(42)

lp <- matrix(rnorm(50, mean = 0, sd = 0.5), nrow = 50)

result <- sim_tte(
  pi       = lp,
  mu       = -1,
  coefs    = 1.1,
  time     = seq(0.1, 100, by = 0.1),
  type     = "weibull",
  end_time = 100
)
head(result)
\end{verbatim}

For subject \(i\), the Weibull hazard is given by \(h_i(t) = \exp(\mu + \text{lp}_i) \cdot \gamma \cdot t^{\gamma - 1}\),
where \(\mu\) is the intercept, \(\text{lp}_i\) is the subject specific log prognostic index, and \(\gamma > 0\) is the shape parameter, supplied through \texttt{coefs}. When \(\gamma = 1\) the model reduces to the exponential distribution with constant hazard; values of \(\gamma > 1\) correspond to an increasing hazard over time, whereas \(\gamma < 1\) implies a decreasing hazard.

The function \texttt{explore\_pi\_tq\_surv()} can be used to examine how the range of prognostic index values translates into differences in survival probability at the median survival time.

\begin{verbatim}
data_pi <- explore_pi_tq_surv(
  pi       = seq(-2, 2, by = 0.25),
  mu       = -1,
  shape    = 1.1,
  type     = "weibull",
  end_time = 100
)

plot(survdiff_tq ~ exp(lp), data = data_pi,
     xlab = "Hazard Ratio", ylab = "Survival difference at median",
     log  = "x", type = "l", lwd = 2)
abline(h = 0, v = 1, lty = 2)
\end{verbatim}

The resulting curve relates log hazard ratio values to differences in survival probability at the population median survival time, providing a more interpretable sense of effect size before specifying a prognostic index distribution.

\subsection{M-spline flexible baseline hazard}\label{m-spline-flexible-baseline-hazard}

Weibull models restrict the baseline hazard to a monotone form and may therefore be unsuitable when the hazard exhibits early peaks, plateaus, or more complex time dependence. To accommodate such settings, \CRANpkg{simtte} supports M spline baseline hazards, in which
\(h_0(t)\) is represented as a nonnegative linear combination of M spline basis functions \citep{ramsay1988}:

\[h_0(t) = \sum_{k=1}^{K} \alpha_k M_k(t), \quad \alpha_k \geq 0.\]

The non-negativity of M-spline basis functions ensures that \(h_0(t) \geq 0\) whenever the coefficient vector is nonnegative, making this basis convenient for hazard modeling. This construction is related to flexible parametric survival models such as those of \citet{royston2002}, although the parameterization differs: here the hazard itself is represented directly through an M spline basis, rather than modeling the log cumulative hazard with restricted cubic splines on the log time scale.

The package includes an example M spline specification in \texttt{ms\_data}:

\begin{verbatim}
data("ms_data")
str(ms_data)
\end{verbatim}

The \texttt{basis} component is a 299 by 6 matrix of cubic M spline basis values evaluated at 299 irregularly spaced time points over a 0 to 7.3 time horizon, generated using \texttt{splines2::mSpline()} \citep{splines2}. The product \texttt{basis\ \%*\%\ coefs} gives the baseline hazard evaluated on this grid. Users are not restricted to this specification; any M spline basis matrix generated with \texttt{splines2::mSpline()} can be supplied to \texttt{sim\_tte()}.

The following example defines a new M spline specification over a 24 unit horizon with a broader unimodal hazard shape, and compares the resulting simulated survival curves with those obtained from the bundled example:

\begin{verbatim}
library(splines2)

time_new  <- seq(0.1, 24, by = 0.1)
knots_new <- quantile(time_new, probs = c(0.15, 0.35, 0.55, 0.75))

basis_new <- mSpline(time_new, knots = knots_new,
                     Boundary.knots = c(0, 24),
                     degree = 2, intercept = TRUE)
coefs_new <- c(0.05, 0.35, 0.55, 0.30, 0.10, 0.04, 0.02)

n      <- 200
lp_new <- matrix(rep(0, n), nrow = n)

sim_new <- sim_tte(
  pi     = lp_new,
  mu     = -0.5,
  basis  = basis_new,
  coefs  = coefs_new,
  time   = time_new,
  type   = "ms"
)
mean(sim_new$sim_status)
\end{verbatim}

The simulation call follows the same structure as in the Weibull example, with the M spline basis supplied through \texttt{basis}, the corresponding coefficients through \texttt{coefs}, and \texttt{type\ =\ "ms"} selecting the hazard model. The output has the same standard four column structure and can be used directly with functions such as \texttt{survival::survfit()} for Kaplan Meier estimation.

\subsection{ODE-coupled PK/PD hazard}\label{pkpd-section}

A more general application of \CRANpkg{simtte} is to define the hazard through a mechanistic ODE model. We illustrate this using a standard pharmacometric structure: a one compartment PK model with first order oral absorption coupled to an indirect response PD model \citep{dayneka1993}, in which drug inhibits production of a mediator \(R\) and the mediator level determines the event hazard.

The model is defined by:

\[\frac{dA_\mathrm{gut}}{dt} = -k_a A_\mathrm{gut},\]
\[\frac{dA_\mathrm{cent}}{dt} = k_a A_\mathrm{gut} - \frac{CL_i}{V_i} A_\mathrm{cent},\]
\[\frac{dR}{dt} = k_\mathrm{in}\bigl(1 - \mathrm{INH}(t)\bigr) - k_\mathrm{out} R,\]
\[\mathrm{INH}(t) = \frac{I_\mathrm{max}\, C(t)}{IC_{50,i} + C(t)},\]
\[h(t) = h_0 \frac{R(t)}{R_0},\]

where \(C(t) = A_\mathrm{cent}(t) / V_i\) is the drug concentration,
\(R_0 = k_\mathrm{in} / k_\mathrm{out}\) is the baseline mediator
steady state, and subscript \(i\) denotes individual-level random
effects on clearance, volume, and potency. Under this specification, the hazard equals \(h_0\) at baseline and decreases as treatment suppresses \(R\).

Rather than solving the pharmacological system first and then evaluating survival in a separate step, the survival function \(S(t)\) is introduced as an additional state variable alongside the PK/PD compartments. The corresponding \CRANpkg{mrgsolve} model specification is:

\begin{verbatim}
library(mrgsolve)

code <- '
$PARAM CL=2, V=20, KA=1, KIN=1, KOUT=0.5, IC50=1, IMAX=0.9, H0=0.15

$CMT GUT CENT R

$INIT p11 = 1

$OMEGA @labels ECL EV EIC50
0.09 0.09 0.16

$MAIN
double CLi   = CL * exp(ECL);
double Vi    = V  * exp(EV);
double IC50i = IC50 * exp(EIC50);
R_0 = KIN / KOUT;

$ODE
double CP  = CENT / Vi;
double INH = IMAX * CP / (IC50i + CP);
double HAZ = H0 * (R / (KIN / KOUT));
dxdt_GUT  = -KA * GUT;
dxdt_CENT =  KA * GUT - (CLi / Vi) * CENT;
dxdt_R    =  KIN * (1 - INH) - KOUT * R;
dxdt_p11  = -p11 * HAZ;

$CAPTURE CP HAZ
'
mod <- mcode("pkpd_hazard", code, quiet = TRUE)
\end{verbatim}

The term \texttt{dxdt\_p11\ =\ -p11\ *\ HAZ} implements the survival component. Here, \texttt{HAZ} is evaluated within the \texttt{\$ODE} block from the current model states, so the solver updates the pharmacological variables and survival probability simultaneously. The full hazard and survival calculation therefore remains within the compiled model code.

Between subject variability is specified in \texttt{\$OMEGA} on log clearance, log volume, and log potency. With the values shown above, the implied coefficients of variation are approximately 30\% for clearance and volume and 42\% for \(IC_{50}\), yielding subject specific exposure, response, hazard, and survival trajectories.

Event times are generated by simulating a cohort under a once daily oral dosing regimen and passing the resulting survival trajectories to \texttt{sim\_tte\_df()}:

\begin{verbatim}
library(dplyr)
library(survival)

set.seed(20260714)
n         <- 300
data_full <- expand.ev(ID = 1:n, amt = 100, cmt = 1, ii = 1,
                       addl = 23, time = 0)
out_full  <- as.data.frame(mrgsim(mod, data = data_full,
                                  end = 24, delta = 0.1))

sim_event <- sim_tte_df(out_full %>% select(ID, time, p11))
mean(sim_event$sim_status)   # event rate before additional censoring
\end{verbatim}

The call to \texttt{sim\_tte\_df()} does not depend on how the survival trajectory was generated; the required input is a data frame containing subject identifier, time, and survival probability.

As a numerical check, one can compare the Kaplan Meier estimate from a censored simulated cohort with the mean \texttt{p11} trajectory across subjects. Under noninformative censoring, the Kaplan Meier estimator targets the marginal survival function, so the two summaries would be expected to be close. This kind of self-consistency check is complemented by package-level validation using a compact PK/PD model whose cumulative hazard has an independent closed form (linear drug effect on a baseline hazard), for which the simulated event-time distribution agreed with the closed-form reference to within statistical tolerance, and the event rate responded monotonically to the strength of the drug effect, as expected mechanistically (package validation scripts, \texttt{inst/validation/}):

\begin{verbatim}
apply_censoring <- function(event_time, true_status, shape, scale) {
  cens <- rweibull(length(event_time), shape = shape, scale = scale)
  data.frame(
    observed_time = pmin(event_time, cens),
    status        = as.integer(true_status == 1 & event_time <= cens)
  )
}

cohort    <- apply_censoring(sim_event$sim_time, sim_event$sim_status,
                              shape = 1.5, scale = 20)
fit_km    <- survfit(Surv(cohort$observed_time, cohort$status) ~ 1)
mean_p11  <- out_full %>% group_by(time) %>%
             summarise(mean_p11 = mean(p11), .groups = "drop")

plot(fit_km, col = "#CC3333", lwd = 2, conf.int = FALSE,
     xlab = "Time (days)", ylab = "Survival probability",
     main = "Simulated KM vs. population-average p11")
lines(mean_p11$time, mean_p11$mean_p11, col = "black", lwd = 2, lty = 2)
legend("topright",
       legend = c("KM (simulated, censored)", "Mean p11 (ODE)"),
       col = c("#CC3333", "black"), lwd = 2, lty = c(1, 2), bty = "n")
\end{verbatim}

To illustrate the effect of mechanistic parameters on the event process, the simulation can be repeated at different values of \(IC_{50}\):

\begin{verbatim}
for (ic50 in c(0.3, 1.0, 5.0)) {
  out_x <- as.data.frame(mrgsim(mod, data = data_full,
                                end = 24, delta = 0.1,
                                param = list(IC50 = ic50)))
  sim_x <- sim_tte_df(out_x %>% select(ID, time, p11))
  cat("IC50 =", ic50, "-> event rate =",
      round(mean(sim_x$sim_status), 3), "\n")
}
# IC50 = 0.3  -> event rate = 0.503
# IC50 = 1.0  -> event rate = 0.560
# IC50 = 5.0  -> event rate = 0.727
\end{verbatim}

In this example, lower values of \(IC_{50}\) correspond to greater potency, greater suppression of the mediator, and a lower event rate over the simulated follow up period. The example illustrates how changes in a pharmacological parameter propagate through the PK/PD model to alter the survival outcome without requiring a separate empirical hazard specification.

\section{Censoring mechanisms}\label{censoring}

\subsection{Administrative censoring}\label{administrative-censoring}

\texttt{sim\_tte()} applies administrative censoring natively through the
\texttt{end\_time} argument, which defaults to the maximum value of the
\texttt{time} vector and is threaded explicitly through the simulation
pipeline. Any subject whose event time under the hazard model
exceeds \texttt{end\_time} receives \texttt{sim\_status\ =\ 0} and
\texttt{sim\_time\ =\ end\_time}. This corresponds to a study with a fixed
analysis cutoff at which all remaining subjects are censored
simultaneously. This behavior is unaffected by \texttt{event\_time\_method}:
interpolation (see Event-time interpolation above) only refines the
time associated with a crossing already detected on the reported
grid, and never alters whether or when a subject is censored.

\subsection{Random censoring}\label{random-censoring}

Independent random censoring is not applied within the simulation functions, but it can be added directly to the simulated event times. For each subject, the observed time is defined as the minimum of the event time and the censoring time, and the event indicator records whether the event occurred before censoring:

\begin{verbatim}
set.seed(1)
# event_time and true_status are from sim_tte() or sim_tte_df()
cens_time     <- rweibull(nrow(result), shape = 1.5, scale = 30)
observed_time <- pmin(result$sim_time, cens_time)
status        <- as.integer(result$sim_time <= cens_time &
                              result$sim_status == 1)
\end{verbatim}

This construction assumes that censoring is statistically independent of the event process, corresponding to the usual noninformative censoring assumption underlying Kaplan Meier estimation, Cox regression, and related survival methods. In simulation settings where informative censoring is of interest, for example when dropout depends on the same subject characteristics that influence event risk, censoring times can instead be generated from a distribution parameterized by individual level quantities such as the prognostic index. In that case, the dependence structure is specified explicitly by the user.

This design leaves censoring outside the core simulation step. Because \texttt{sim\_tte()} and \texttt{sim\_tte\_df()} return standard data frames containing simulated event times and event indicators, censoring can be incorporated using ordinary R operations without requiring additional package specific machinery.

\section{Summary and discussion}\label{discussion}

\subsection{Summary}\label{summary}

\CRANpkg{simtte} provides an R interface for simulating time to event data under three model classes: Weibull, M spline, and user defined ODE based hazards, with a common interface and output structure across these settings. Its main methodological feature is the joint treatment of survival and pharmacological dynamics within the same ODE system. By adding a survival compartment through (\texttt{dxdt\_p11\ =\ -p11\ *\ HAZ}) to an \CRANpkg{mrgsolve} model, the hazard can be defined as a function of arbitrary time varying model states. The resulting survival trajectory is obtained directly from the numerical solution, and event times are then generated by inverse transform sampling.

This formulation extends time to event simulation beyond closed form hazard specifications and supports a more direct connection between mechanistic PK/PD models and simulated clinical endpoints. In settings where event risk is expected to emerge from exposure, biomarker dynamics, or disease progression, it provides a practical way to represent that dependence within an established pharmacometric workflow. The implications, current limitations, and possible methodological extensions are discussed in the following sections.

\subsection{Limitations}\label{limitations}

A primary limitation of the current implementation is that, under the default \texttt{event\_time\_method\ =\ "grid"}, event time resolution is determined by the reported output grid, which is controlled by the \texttt{time} argument (Section \ref{package-design}) rather than by \CRANpkg{mrgsolve}'s own default schedule. If \texttt{time} is coarse, the resulting event times will be correspondingly discretized. Setting \texttt{event\_time\_method\ =\ "log\_survival"} mitigates this for a crossing that falls strictly between two reported points, by linearly interpolating cumulative hazard, but this recovers the exact event time only when the hazard is genuinely piecewise constant on the reported grid, as for the M spline model; for the Weibull model with \(\gamma \neq 1\), or for a custom mechanistic model whose hazard varies continuously between reported points, it remains an approximation whose accuracy depends on the coarseness of \texttt{time}. Applications requiring finer temporal precision therefore still benefit from a finer \texttt{time} grid, with the usual trade off in computational cost; for the built in Weibull model this cost scales roughly linearly with cohort size and only weakly with grid density (package validation scripts, \texttt{inst/validation/}).

The package currently supports simulation of a single event type per subject. In principle, competing or multiple event types could be handled by augmenting the ODE model with one survival compartment per event process and applying the event time generation step separately to each trajectory, followed by comparison of the resulting event times. However, this approach is not yet formalized in the current package interface.

\subsection{Future directions}\label{future-directions}

Several extensions follow naturally from the current framework. One priority is support for competing risks, recurrent events, and multi state event processes, particularly in cardiovascular, renal, inflammatory, and other disease areas applications where repeated and heterogeneous event types are common. COPD exacerbation studies provide one example in which recurrent events and distinct clinical event processes may both be relevant. The survival co integration approach can be extended by assigning a separate event process to each cause specific or recurrent event hazard, with each hazard defined as a function of the underlying ODE states. A higher level interface for these settings would reduce the manual coordination currently required across multiple event processes.

A second direction is the propagation of parameter uncertainty from PK/PD models to trial level operating characteristics. The current implementation generates event times conditional on a fixed set of model parameters. An extension could instead sample from a parameter distribution, for example a posterior distribution from a Bayesian population PK analysis or an estimated uncertainty distribution from a nonlinear mixed effects model, and propagate that uncertainty to quantities such as event rates, survival probabilities, or power. This would provide a clearer connection between model based uncertainty and simulation based trial evaluation.

A third area for development is more explicit support for censoring mechanisms beyond independent random censoring. In particular, informative censoring could be represented by allowing the censoring process to depend on the same ODE states or subject level factors that influence event risk. For example, a censoring hazard linked to drug exposure, biomarker dynamics, or disease status would allow investigation of the robustness of standard survival methods under mechanistically motivated dropout scenarios.

\bibliography{RJreferences.bib}

\address{%
Carlos Serra Traynor\\
Clinical Pharmacology \& Quantitative Pharmacology, Clinical Pharmacology \& Safety Sciences, R\&D BioPharmaceuticals, AstraZeneca\\%
Barcelona, Spain\\
%
%
%
\href{mailto:carlos.traynor.qcp@gmail.com}{\nolinkurl{carlos.traynor.qcp@gmail.com}}%
}
